\scanpage{7}
\mspage{5}
\[
\begin{aligned}
\underline{\mathrm{Ch}}(X)&\longrightarrow X\times X\\
c&\longmapsto(\sigma(x),\beta(\unsure{x|m}))\quad.
\end{aligned}
\]
Pour tout chemin $c$ dans $X$ i.e. un diagramme $(*)$ (satisfaisant à la condition de dégénérescence à l'infini), on a une source ou origine du chemin, et un but ou une \uline{extrémité}
\[
\left\{\begin{aligned}
\sigma(x)&=A_i\quad\text{pour }\deleted{|i|\ \text{grand}}\ \deleted{i\leq-N,\ N\ \text{grand}\ \unsure{(-N}}\,i\leq-\operatorname{long}(x)\\
\deleted{\unsure{\sigma}}\beta(x)&=A_i\quad\text{pour }i\geq\operatorname{long}(\unsure{x|a})
\end{aligned}\right.
\]
La fibre du \deleted{\illegible{1}} foncteur précédent en $x,y$ sera notée
\[
\underline{\mathrm{Ch}}(X;x,y)\quad\text{ou simplement }\underline{\mathrm{Ch}}(x,y)\quad,
\]
\uline{catégorie des chemins de} $x$ à $y$.

2. \deleted{\unsure{Le}} \uline{Calcul d'une cat. de fractions} $X\Sigma^{-1}$

Avant de continuer le formalisme des chemins (composition etc), je vais \unsure{tester|tâter} l'utilité de la notion de \deleted{\unsure{deux}} catégorie de chemins, par le ``calcul'' des $X\Sigma^{-1}$, où
\[
\Sigma\subset\operatorname{Fl}(X)\qquad\inserted{\Sigma\ \text{contient les identités}}
\]
est une partie \inserted{\unsure{autrement|contenant les $1$}} \uline{quelconque} de $\operatorname{Fl}(X)$. On sait que $\operatorname{Ob}\deleted{\unsure{F}}X\Sigma^{-1}=\operatorname{Ob}X$, et tout revient à définir \unsure{se|s}
\deleted{\unsure{pr}} les
\[
\operatorname{Hom}_{X\Sigma^{-1}}(x,y)\qquad x,y\in\operatorname{Ob}X,
\]

% archive sheet number: 6
